\documentclass[12pt,a4paper]{exam}

% --------------------
% Packages
% --------------------
\usepackage[utf8]{inputenc}
\usepackage[T1]{fontenc}
\usepackage[french]{babel}
\usepackage{amsmath,amssymb}
\usepackage{geometry}
\usepackage{pgfplots}
\pgfplotsset{compat=1.18}
\geometry{margin=2cm}

% --------------------
% En-tête / pied (exam)
% --------------------
\header{}{DS de mathématiques}{1ère STMG}
\footer{}{\thepage}{}

% --------------------
\begin{document}

\begin{center}
    \Large\textbf{Devoir surveillé\\Fonctions polynômes de degré 2}
\end{center}

\vspace{0.3cm}

\noindent
\textbf{Durée :} 1h \hfill
\textbf{Calculatrice :} \underline{NON} autorisée

\vspace{0.5cm}

% ======================================================
\noindent
\textbf{Exercice 1 : Étude d’une fonction (\ \ \ \  / 5 points )}
\vspace{5pt}
\noindent\\On représente ci-dessous la courbe de la fonction $f(x) = ax^2+c$.

\begin{center}
    \begin{tikzpicture}[scale=0.8]

    % Grille
    \draw[step=1cm, gray, very thin] (-3,-5) grid (3,5);

    % Axes
    \draw[->] (-3.2,0) -- (3.2,0) node[right] {$x$};
    \draw[->] (0,-5.2) -- (0,5.2) node[above] {$y$};

    % Graduations x
    \foreach \x in {-3,-2,-1,1,2,3}
        \draw (\x,0.1) -- (\x,-0.1) node[below] {\small $\x$};

    % Graduations y
    \foreach \y in {-4,-2,2,4}
        \draw (0.1,\y) -- (-0.1,\y) node[left] {\small $\y$};

    % Courbe y = 3x^2 - 4 (limitée à y ≤ 5)
    \draw[domain=-sqrt(2.7):sqrt(2.7), samples=200, smooth, thick, red]
        plot (\x, {3*\x*\x - 3});

    \end{tikzpicture}
\end{center}


\begin{questions}

\question[2] Déterminer graphiquement les valeurs de $f(0)$ et $f(1)$.

\question[3] Retrouver les valeurs de $a$ et $c$ à l'aide de la question
précédente.

\end{questions}

\vspace{0.5cm}

% ======================================================
\noindent
\textbf{Exercice 2 : Étudier et tracer ( \ \ \ / 8 points )}

\vspace{5pt}
\noindent
On considère la fonction définie sur $\mathbb{R}$ suivante :
\[
f(x) = -2(x+1)(x-5)
\]

\begin{questions}

\question Trouver les caractéristiques suivantes du polynôme 
$-2(x+1)(x-5)$ :
\begin{parts}
    \part[2] Racines
    \part[1] Équation de l'axe de symétrie 
    \part[1] Coordonnées du sommet 
    \part[1] Le sommet sera-t-il un maximum ou un minimum ? 
    \textbf{\underline{Justifier}}
\end{parts}

\question[1] Tracer l'allure de la courbe représentative de $f$.

% \question Construire le tableau de signes de $f$.

\question[2] Construire le tableau de \underline{variations} de $f$. On fera 
apparaître les valeurs au sommet.

\end{questions}

\vspace{0.5cm}

% ======================================================
\noindent
\textbf{Exercice 3 : Problème ( \ \ \ / 7 points )}

\vspace{5pt}
\noindent
Pour tout réel $x$, on pose $f(x) = -2x^2 + 2x + 4$.

\begin{questions}

\question[1] Montrer que $2$ est une racine du 
polynôme $-2x^2 + 2x + 4$.

\question[1] Calculer $f(-1)$. Que peut-on en déduire ?

\question[2] En déduire une forme factorisée de $f(x)$.

\question[2] Construire le tableau de signes de $f$ sur $\mathbb{R}$.

\question[1] Résoudre l'inéquation $f(x) \leq 0$ sur $\mathbb{R}$.

\end{questions}

\vspace{0.5cm}

\begin{center}
    \textit{Fin du devoir}
\end{center}

\end{document}
